% Options for packages loaded elsewhere
% Options for packages loaded elsewhere
\PassOptionsToPackage{unicode}{hyperref}
\PassOptionsToPackage{hyphens}{url}
%
\documentclass[
  ignorenonframetext,
  aspectratio=169,
]{beamer}
\newif\ifbibliography
\usepackage{pgfpages}
\setbeamertemplate{caption}[numbered]
\setbeamertemplate{caption label separator}{: }
\setbeamercolor{caption name}{fg=normal text.fg}
\beamertemplatenavigationsymbolsempty
% remove section numbering
\setbeamertemplate{part page}{
  \centering
  \begin{beamercolorbox}[sep=16pt,center]{part title}
    \usebeamerfont{part title}\insertpart\par
  \end{beamercolorbox}
}
\setbeamertemplate{section page}{
  \centering
  \begin{beamercolorbox}[sep=12pt,center]{section title}
    \usebeamerfont{section title}\insertsection\par
  \end{beamercolorbox}
}
\setbeamertemplate{subsection page}{
  \centering
  \begin{beamercolorbox}[sep=8pt,center]{subsection title}
    \usebeamerfont{subsection title}\insertsubsection\par
  \end{beamercolorbox}
}
% Prevent slide breaks in the middle of a paragraph
\widowpenalties 1 10000
\raggedbottom
\AtBeginPart{
  \frame{\partpage}
}
\AtBeginSection{
  \ifbibliography
  \else
    \frame{\sectionpage}
  \fi
}
\AtBeginSubsection{
  \frame{\subsectionpage}
}
\usepackage{iftex}
\ifPDFTeX
  \usepackage[T1]{fontenc}
  \usepackage[utf8]{inputenc}
  \usepackage{textcomp} % provide euro and other symbols
\else % if luatex or xetex
  \usepackage{unicode-math} % this also loads fontspec
  \defaultfontfeatures{Scale=MatchLowercase}
  \defaultfontfeatures[\rmfamily]{Ligatures=TeX,Scale=1}
\fi
\usepackage{lmodern}

\usetheme[]{Boadilla}
\usecolortheme[]{default}
\usefonttheme[]{professionalfonts}
\ifPDFTeX\else
  % xetex/luatex font selection
\fi
% Use upquote if available, for straight quotes in verbatim environments
\IfFileExists{upquote.sty}{\usepackage{upquote}}{}
\IfFileExists{microtype.sty}{% use microtype if available
  \usepackage[]{microtype}
  \UseMicrotypeSet[protrusion]{basicmath} % disable protrusion for tt fonts
}{}
\makeatletter
\@ifundefined{KOMAClassName}{% if non-KOMA class
  \IfFileExists{parskip.sty}{%
    \usepackage{parskip}
  }{% else
    \setlength{\parindent}{0pt}
    \setlength{\parskip}{6pt plus 2pt minus 1pt}}
}{% if KOMA class
  \KOMAoptions{parskip=half}}
\makeatother

\usepackage{color}
\usepackage{fancyvrb}
\newcommand{\VerbBar}{|}
\newcommand{\VERB}{\Verb[commandchars=\\\{\}]}
\DefineVerbatimEnvironment{Highlighting}{Verbatim}{commandchars=\\\{\}}
% Add ',fontsize=\small' for more characters per line
\usepackage{framed}
\definecolor{shadecolor}{RGB}{241,243,245}
\newenvironment{Shaded}{\begin{snugshade}}{\end{snugshade}}
\newcommand{\AlertTok}[1]{\textcolor[rgb]{0.68,0.00,0.00}{#1}}
\newcommand{\AnnotationTok}[1]{\textcolor[rgb]{0.37,0.37,0.37}{#1}}
\newcommand{\AttributeTok}[1]{\textcolor[rgb]{0.40,0.45,0.13}{#1}}
\newcommand{\BaseNTok}[1]{\textcolor[rgb]{0.68,0.00,0.00}{#1}}
\newcommand{\BuiltInTok}[1]{\textcolor[rgb]{0.00,0.23,0.31}{#1}}
\newcommand{\CharTok}[1]{\textcolor[rgb]{0.13,0.47,0.30}{#1}}
\newcommand{\CommentTok}[1]{\textcolor[rgb]{0.37,0.37,0.37}{#1}}
\newcommand{\CommentVarTok}[1]{\textcolor[rgb]{0.37,0.37,0.37}{\textit{#1}}}
\newcommand{\ConstantTok}[1]{\textcolor[rgb]{0.56,0.35,0.01}{#1}}
\newcommand{\ControlFlowTok}[1]{\textcolor[rgb]{0.00,0.23,0.31}{\textbf{#1}}}
\newcommand{\DataTypeTok}[1]{\textcolor[rgb]{0.68,0.00,0.00}{#1}}
\newcommand{\DecValTok}[1]{\textcolor[rgb]{0.68,0.00,0.00}{#1}}
\newcommand{\DocumentationTok}[1]{\textcolor[rgb]{0.37,0.37,0.37}{\textit{#1}}}
\newcommand{\ErrorTok}[1]{\textcolor[rgb]{0.68,0.00,0.00}{#1}}
\newcommand{\ExtensionTok}[1]{\textcolor[rgb]{0.00,0.23,0.31}{#1}}
\newcommand{\FloatTok}[1]{\textcolor[rgb]{0.68,0.00,0.00}{#1}}
\newcommand{\FunctionTok}[1]{\textcolor[rgb]{0.28,0.35,0.67}{#1}}
\newcommand{\ImportTok}[1]{\textcolor[rgb]{0.00,0.46,0.62}{#1}}
\newcommand{\InformationTok}[1]{\textcolor[rgb]{0.37,0.37,0.37}{#1}}
\newcommand{\KeywordTok}[1]{\textcolor[rgb]{0.00,0.23,0.31}{\textbf{#1}}}
\newcommand{\NormalTok}[1]{\textcolor[rgb]{0.00,0.23,0.31}{#1}}
\newcommand{\OperatorTok}[1]{\textcolor[rgb]{0.37,0.37,0.37}{#1}}
\newcommand{\OtherTok}[1]{\textcolor[rgb]{0.00,0.23,0.31}{#1}}
\newcommand{\PreprocessorTok}[1]{\textcolor[rgb]{0.68,0.00,0.00}{#1}}
\newcommand{\RegionMarkerTok}[1]{\textcolor[rgb]{0.00,0.23,0.31}{#1}}
\newcommand{\SpecialCharTok}[1]{\textcolor[rgb]{0.37,0.37,0.37}{#1}}
\newcommand{\SpecialStringTok}[1]{\textcolor[rgb]{0.13,0.47,0.30}{#1}}
\newcommand{\StringTok}[1]{\textcolor[rgb]{0.13,0.47,0.30}{#1}}
\newcommand{\VariableTok}[1]{\textcolor[rgb]{0.07,0.07,0.07}{#1}}
\newcommand{\VerbatimStringTok}[1]{\textcolor[rgb]{0.13,0.47,0.30}{#1}}
\newcommand{\WarningTok}[1]{\textcolor[rgb]{0.37,0.37,0.37}{\textit{#1}}}

\usepackage{longtable,booktabs,array}
\newcounter{none} % for unnumbered tables
\usepackage{calc} % for calculating minipage widths
\usepackage{caption}
% Make caption package work with longtable
\makeatletter
\def\fnum@table{\tablename~\thetable}
\makeatother
\usepackage{graphicx}
\makeatletter
\newsavebox\pandoc@box
\newcommand*\pandocbounded[1]{% scales image to fit in text height/width
  \sbox\pandoc@box{#1}%
  \Gscale@div\@tempa{\textheight}{\dimexpr\ht\pandoc@box+\dp\pandoc@box\relax}%
  \Gscale@div\@tempb{\linewidth}{\wd\pandoc@box}%
  \ifdim\@tempb\p@<\@tempa\p@\let\@tempa\@tempb\fi% select the smaller of both
  \ifdim\@tempa\p@<\p@\scalebox{\@tempa}{\usebox\pandoc@box}%
  \else\usebox{\pandoc@box}%
  \fi%
}
% Set default figure placement to htbp
\def\fps@figure{htbp}
\makeatother





\setlength{\emergencystretch}{3em} % prevent overfull lines

\providecommand{\tightlist}{%
  \setlength{\itemsep}{0pt}\setlength{\parskip}{0pt}}



 


\makeatletter
\@ifpackageloaded{caption}{}{\usepackage{caption}}
\AtBeginDocument{%
\ifdefined\contentsname
  \renewcommand*\contentsname{Table of contents}
\else
  \newcommand\contentsname{Table of contents}
\fi
\ifdefined\listfigurename
  \renewcommand*\listfigurename{List of Figures}
\else
  \newcommand\listfigurename{List of Figures}
\fi
\ifdefined\listtablename
  \renewcommand*\listtablename{List of Tables}
\else
  \newcommand\listtablename{List of Tables}
\fi
\ifdefined\figurename
  \renewcommand*\figurename{Figure}
\else
  \newcommand\figurename{Figure}
\fi
\ifdefined\tablename
  \renewcommand*\tablename{Table}
\else
  \newcommand\tablename{Table}
\fi
}
\@ifpackageloaded{float}{}{\usepackage{float}}
\floatstyle{ruled}
\@ifundefined{c@chapter}{\newfloat{codelisting}{h}{lop}}{\newfloat{codelisting}{h}{lop}[chapter]}
\floatname{codelisting}{Listing}
\newcommand*\listoflistings{\listof{codelisting}{List of Listings}}
\makeatother
\makeatletter
\makeatother
\makeatletter
\@ifpackageloaded{caption}{}{\usepackage{caption}}
\@ifpackageloaded{subcaption}{}{\usepackage{subcaption}}
\makeatother

\usepackage{bookmark}
\IfFileExists{xurl.sty}{\usepackage{xurl}}{} % add URL line breaks if available
\urlstyle{same}
\hypersetup{
  pdftitle={Chapter 3: Risk Models},
  hidelinks,
  pdfcreator={LaTeX via pandoc}}


\title{\texorpdfstring{Chapter 3: Risk Models}{Chapter 3: Risk Models}}
\author{}
\date{}

\begin{document}
\frame{\titlepage}


\section{Compound Binomial Distributions}\label{compound-binomial}

\begin{frame}{From Individual Policies to Aggregate Loss}
\protect\phantomsection\label{from-individual-policies-to-aggregate-loss}
Consider a portfolio of \(n\) independent and identical policies.

Each policy can give rise to \textbf{at most one claim} during the
period.

Let \(p\) be the probability that a policy produces a claim.

The number of claims in the portfolio is therefore

\[
N \sim \operatorname{Binomial}(n,p).
\]
\end{frame}

\begin{frame}{Why the Binomial Distribution?}
\protect\phantomsection\label{why-the-binomial-distribution}
The binomial model is natural when:

\begin{itemize}
\tightlist
\item
  there are a fixed number \(n\) of potential claim events;
\item
  each policy produces either \textbf{zero or one claim};
\item
  each policy has the same claim probability \(p\);
\item
  claim events are independent.
\end{itemize}

Thus, the number of claims has an upper limit:

\[
0 \leq N \leq n.
\]
\end{frame}

\begin{frame}{The Compound Binomial Model}
\protect\phantomsection\label{the-compound-binomial-model}
Let \(X_i\) denote the amount of the \(i\)th claim.

Assume that \(X_1,X_2,\ldots\) are i.i.d. and independent of \(N\).

Then the aggregate claim amount is

\[
S=\sum_{i=1}^{N}X_i.
\]

We call \(S\) a \textbf{compound binomial random variable}.

The distribution is denoted by

\[
S\sim\mathcal{CB}(n,p,F_X).
\]
\end{frame}

\begin{frame}{Aggregate Loss Model}
\protect\phantomsection\label{aggregate-loss-model}
The collective risk model is

\[
S=\sum_{i=1}^{N}X_i,
\]

where

\[
N\sim\operatorname{Binomial}(n,p).
\]

Here, \(N\) determines \textbf{how many claims occur}, while \(X_i\)
determines \textbf{how much each claim costs}.

Thus, the aggregate loss combines a \textbf{frequency model} and a
\textbf{severity model}.
\end{frame}

\begin{frame}{Expected Aggregate Loss}
\protect\phantomsection\label{expected-aggregate-loss}
For a random sum,

\[
S=\sum_{i=1}^{N}X_i,
\]

the expected aggregate loss is

\[
\boxed{E[S]=E[N]E[X]}.
\]

This follows from the standard result for a random sum when \(N\) is
independent of the claim amounts.
\end{frame}

\begin{frame}{Binomial Frequency Model}
\protect\phantomsection\label{binomial-frequency-model}
For the binomial frequency model,

\[
N\sim\operatorname{Binomial}(n,p),
\]

so

\[
E[N]=np.
\]

Let

\[
m_1=E[X].
\]

Substituting into the random-sum formula,

\[
\begin{aligned}
E[S]
&=E[N]E[X]\\
&=(np)m_1.
\end{aligned}
\]

Therefore,

\[
\boxed{E[S]=npm_1}.
\]
\end{frame}

\begin{frame}{Interpretation of Expected Aggregate Loss}
\protect\phantomsection\label{interpretation-of-expected-aggregate-loss}
The formula

\[
E[S]=npm_1
\]

can be written as

\[
\boxed{
E[S]
=
\underbrace{np}_{\text{expected number of claims}}
\underbrace{m_1}_{\text{expected claim amount}}
}.
\]

Thus, expected aggregate loss is determined by:

\begin{itemize}
\tightlist
\item
  the expected number of claims, \(np\);
\item
  the expected amount per claim, \(m_1\).
\end{itemize}

The binomial frequency model therefore gives a simple interpretation of
the expected aggregate loss.
\end{frame}

\begin{frame}{Variance of Aggregate Loss}
\protect\phantomsection\label{variance-of-aggregate-loss}
For a random sum,

\[
\operatorname{Var}(S)
=
E[N]\operatorname{Var}(X)
+
\operatorname{Var}(N)\{E[X]\}^2.
\]

For

\[
N\sim\operatorname{Binomial}(n,p),
\]

we have

\[
E[N]=np,
\qquad
\operatorname{Var}(N)=np(1-p).
\]

Let

\[
m_1=E[X],
\qquad
m_2=E[X^2].
\]
\end{frame}

\begin{frame}{Substituting the Binomial Moments}
\protect\phantomsection\label{substituting-the-binomial-moments}
Since

\[
\operatorname{Var}(X)=m_2-m_1^2,
\]

substitute into the random-sum formula:

\[
\begin{aligned}
\operatorname{Var}(S)
&=np(m_2-m_1^2)
+np(1-p)m_1^2.
\end{aligned}
\]

Expanding,

\[
\begin{aligned}
\operatorname{Var}(S)
&=npm_2-npm_1^2
+np(1-p)m_1^2.
\end{aligned}
\]
\end{frame}

\begin{frame}{Simplifying the Variance}
\protect\phantomsection\label{simplifying-the-variance}
Collect the terms involving \(m_1^2\):

\[
\begin{aligned}
\operatorname{Var}(S)
&=npm_2
-npm_1^2
+np(1-p)m_1^2\\
&=npm_2-np^2m_1^2.
\end{aligned}
\]

Therefore,

\[
\boxed{
\operatorname{Var}(S)
=
npm_2-np^2m_1^2
}.
\]

The result combines \textbf{claim-size uncertainty} and
\textbf{frequency uncertainty} under the binomial model.

I would split the setup and the questions so that the first slide
introduces the model and the second clearly states what students need to
derive.
\end{frame}

\begin{frame}{Worked Example: Compound Binomial with Gamma Severity}
\protect\phantomsection\label{worked-example-compound-binomial-with-gamma-severity}
Suppose

\[
N\sim\operatorname{Binomial}(100,0.01),
\]

and, independently,

\[
X_i\overset{\text{i.i.d.}}{\sim}\operatorname{Gamma}(20,0.1),
\]

where the Gamma distribution uses the \textbf{shape-rate}
parameterisation.

Let

\[
S=\sum_{i=1}^{N}X_i.
\]
\end{frame}

\begin{frame}{Worked Example: Tasks}
\protect\phantomsection\label{worked-example-tasks}
For the compound binomial model above, find:

\begin{enumerate}
\tightlist
\item
  the MGF of \(S\);
\item
  \(E[S]\) using the MGF;
\item
  \(\operatorname{Var}(S)\) using the MGF;
\item
  \(P(S=0)\).
\end{enumerate}
\end{frame}

\begin{frame}{1. MGF of the Aggregate Loss}
\protect\phantomsection\label{mgf-of-the-aggregate-loss}
For

\[
X\sim\operatorname{Gamma}(20,0.1),
\]

the MGF is

\[
M_X(t)
=
\left(\frac{0.1}{0.1-t}\right)^{20}.
\]

For the compound binomial model,

\[
M_S(t)
=
\left[q+pM_X(t)\right]^n.
\]

Here, \(n=100,\qquad p=0.01,\qquad q=0.99.\)
\end{frame}

\begin{frame}{MGF of \(S\)}
\protect\phantomsection\label{mgf-of-s}
Substituting the values,

\[
\boxed{
M_S(t)=
\left[
0.99+
0.01\left(\frac{0.1}{0.1-t}\right)^{20}
\right]^{100}
}
\]

for \(t<0.1\).

This follows directly from the compound binomial MGF.
\end{frame}

\begin{frame}{2. Expected Aggregate Loss}
\protect\phantomsection\label{expected-aggregate-loss-1}
Use \(E[S]=M_S'(0).\)

Let

\[
A(t)=
0.99+
0.01\left(\frac{0.1}{0.1-t}\right)^{20}.
\]

Then

\[
M_S(t)=A(t)^{100}.
\]

Hence, \(M_S'(t)=100A(t)^{99}A'(t).\)

At \(t=0\), \qquad \(A(0)=1, \qquad A'(0)=0.01M_X'(0).\)
\end{frame}

\begin{frame}{Expected Aggregate Loss}
\protect\phantomsection\label{expected-aggregate-loss-2}
For the Gamma distribution,

\[
M_X'(0)=E[X]
=\frac{20}{0.1}=200.
\]

Therefore,

\[
\begin{aligned}
E[S]
&=M_S'(0)\\
&=100(0.01)(200).
\end{aligned}
\]

Thus,

\[
\boxed{E[S]=200}.
\]
\end{frame}

\begin{frame}{3. Variance Using the MGF}
\protect\phantomsection\label{variance-using-the-mgf}
Use

\[
\operatorname{Var}(S)
=
M_S''(0)-[M_S'(0)]^2.
\]

From the previous result,

\[
M_S'(0)=200.
\]

We therefore need

\[
M_S''(0).
\]
\end{frame}

\begin{frame}{Second Derivative}
\protect\phantomsection\label{second-derivative}
For

\[
M_S(t)=A(t)^{100},
\]

the second derivative is

\[
M_S''(t)
=
100(99)A(t)^{98}[A'(t)]^2
+
100A(t)^{99}A''(t).
\]

At \(t=0\),

\[
M_S''(0)
=
100(99)[A'(0)]^2
+
100A''(0).
\]
\end{frame}

\begin{frame}{Evaluate the Components}
\protect\phantomsection\label{evaluate-the-components}
For the Gamma severity,

\[
E[X^2]
=
\operatorname{Var}(X)+[E(X)]^2.
\]

Since

\[
\operatorname{Var}(X)=\frac{20}{0.1^2}=2{,}000,
\]

we obtain

\[
E[X^2]=2{,}000+200^2=42{,}000.
\]

Therefore, \(A'(0)=0.01(200), \qquad A''(0)=0.01(42{,}000).\)
\end{frame}

\begin{frame}{Variance of \(S\)}
\protect\phantomsection\label{variance-of-s}
Thus,

\[
M_S''(0)
=
100(99)(0.01\times200)^2
+
100(0.01\times42{,}000).
\]

Hence,

\[
M_S''(0)=81{,}600.
\]

Finally,

\[
\begin{aligned}
\operatorname{Var}(S)
&=81{,}600-200^2\\
&=\boxed{41{,}600}.
\end{aligned}
\]
\end{frame}

\begin{frame}{4. Probability of Zero Aggregate Loss}
\protect\phantomsection\label{probability-of-zero-aggregate-loss}
Because claim amounts are positive,

\[
S=0
\quad\Longleftrightarrow\quad
N=0.
\]

Therefore,

\[
P(S=0)=P(N=0).
\]

For

\[
N\sim\operatorname{Binomial}(100,0.01),
\]

we have

\[
P(N=0)
=
(1-0.01)^{100}.
\]
\end{frame}

\begin{frame}{Evaluating \(P(S=0)\)}
\protect\phantomsection\label{evaluating-ps0}
Thus,

\[
P(S=0)
=
(0.99)^{100}.
\]

Numerically,

\[
\boxed{P(S=0)\approx0.3660}.
\]

So there is approximately a \textbf{36.6\% probability of no aggregate
claims}.
\end{frame}

\begin{frame}{An Alternative MGF View}
\protect\phantomsection\label{an-alternative-mgf-view}
Recall that \(M_S(t)=E[e^{tS}].\)

As \(t\to-\infty\),

\[
e^{tS}\to
\begin{cases}
1, & S=0,\\
0, & S>0.
\end{cases}
\]

Therefore,

\[
P(S=0)
=
\lim_{t\to-\infty}M_S(t).
\]

For this model,
\(\qquad \lim_{t\to-\infty}M_S(t)=(0.99)^{100} \approx0.3660.\)

This provides another useful interpretation of the MGF.
\end{frame}

\begin{frame}{Key Lessons}
\protect\phantomsection\label{key-lessons}
The compound binomial MGF is obtained by substituting the severity MGF
into

\[
M_S(t)=\left[q+pM_X(t)\right]^n.
\]

For this example,
\(\boxed{M_S(t)=\left[0.99+0.01\left(\frac{0.1}{0.1-t}\right)^{20}\right]^{100}}.\)

The MGF can be used to obtain moments:

\[
E[S]=200,
\qquad
\operatorname{Var}(S)=41{,}600.
\]

The model has a positive probability of zero aggregate loss:
\(P(S=0)=(0.99)^{100}\approx0.3660.\)

\textbf{Main idea:} the compound-binomial MGF provides a single
framework for obtaining both \textbf{aggregate-loss moments} and other
distributional information.
\end{frame}

\begin{frame}{R Simulation}
\protect\phantomsection\label{r-simulation}
We can simulate the collective risk model in two stages.

\begin{enumerate}
\item
  \textbf{Simulate the number of claims} \(N\).
\item
  \textbf{Simulate the claim amounts} \(X_1,\ldots,X_N\) and calculate
  the aggregate loss

  \[
  S=\sum_{i=1}^{N}X_i.
  \]
\end{enumerate}

For this example,

\[
N\sim\operatorname{Binomial}(100,0.01)
\]

and

\[
X_i\overset{\text{i.i.d.}}{\sim}\operatorname{Gamma}(20,0.1).
\]
\end{frame}

\begin{frame}[fragile]{Stage 1: Simulate the Number of Claims}
\protect\phantomsection\label{stage-1-simulate-the-number-of-claims}
For each simulated year, generate

\[
N\sim\operatorname{Binomial}(100,0.01).
\]

Each element of \texttt{N} represents the \textbf{number of claims in
one simulated year}.

For example,

\begin{verbatim}
[1] 0 2 1 2 3 0
\end{verbatim}

The simulation produces \(B=100{,}000\) independent insurance years.
\end{frame}

\begin{frame}[fragile]{Stage 2: Simulate Aggregate Loss}
\protect\phantomsection\label{stage-2-simulate-aggregate-loss}
For each simulated year:

\begin{itemize}
\tightlist
\item
  If \(N=0\), then \(S=0\).
\item
  If \(N=n>0\), simulate \(n\) claim amounts.
\item
  Add the claim amounts to obtain \(S\).
\end{itemize}

Thus, each element of \texttt{S} represents one simulated value of

\[
S=\sum_{i=1}^{N}X_i.
\]
\end{frame}

\begin{frame}[fragile]{R Simulation}
\protect\phantomsection\label{r-simulation-1}
We can simulate the aggregate loss directly in R.

\begin{Shaded}
\begin{Highlighting}[]
\CommentTok{\# Simulate aggregate losses}
\FunctionTok{set.seed}\NormalTok{(}\DecValTok{123}\NormalTok{)}
\NormalTok{B }\OtherTok{\textless{}{-}} \DecValTok{100000}
\NormalTok{N }\OtherTok{\textless{}{-}} \FunctionTok{rbinom}\NormalTok{(B, }\AttributeTok{size =} \DecValTok{100}\NormalTok{, }\AttributeTok{prob =} \FloatTok{0.01}\NormalTok{)}
\NormalTok{S }\OtherTok{\textless{}{-}} \FunctionTok{sapply}\NormalTok{(N, }\ControlFlowTok{function}\NormalTok{(n) \{}
  \FunctionTok{sum}\NormalTok{(}\FunctionTok{rgamma}\NormalTok{(n, }\AttributeTok{shape =} \DecValTok{20}\NormalTok{, }\AttributeTok{rate =} \FloatTok{0.1}\NormalTok{))}
\NormalTok{\})}

\CommentTok{\# Plot the simulated aggregate losses}
\FunctionTok{hist}\NormalTok{(S, }\AttributeTok{breaks =} \DecValTok{50}\NormalTok{,}
     \AttributeTok{main =} \StringTok{"Simulated Aggregate Loss"}\NormalTok{,}
     \AttributeTok{xlab =} \StringTok{"Aggregate loss, S"}\NormalTok{)}
\end{Highlighting}
\end{Shaded}
\end{frame}

\begin{frame}{R Simulation}
\protect\phantomsection\label{r-simulation-2}
We can simulate the aggregate loss directly in R.

\begin{center}
\pandocbounded{\includegraphics[keepaspectratio]{temp-cb_files/figure-beamer/unnamed-chunk-4-1.pdf}}
\end{center}
\end{frame}

\begin{frame}[fragile]{Simulation Results}
\protect\phantomsection\label{simulation-results}
Use the simulated aggregate losses to estimate the moments.

\begin{verbatim}
[1] 199.89
\end{verbatim}

\begin{verbatim}
[1] 41515.39
\end{verbatim}

For \(100{,}000\) simulations, \(\widehat{E[S]}=199.89\) and
\(\widehat{\operatorname{Var}(S)}=41{,}515.39\).

Compare with the theoretical results:

\[
E[S]=200,
\qquad
\operatorname{Var}(S)=41{,}600.
\]

\textbf{The simulation closely reproduces the theoretical results.}
\end{frame}

\begin{frame}{From Model to Simulation}
\protect\phantomsection\label{from-model-to-simulation}
The simulation follows exactly the structure of the collective risk
model:

\[
\boxed{
N
\longrightarrow
X_1,\ldots,X_N
\longrightarrow
S
}
\] \textbf{Frequency:} \(N\sim\operatorname{Binomial}(100,0.01)\)

\textbf{Severity:} \(X_i\sim\operatorname{Gamma}(20,0.1)\)

\textbf{Aggregate loss:} \(S=\sum_{i=1}^{N}X_i\)

\begin{quote}
\textbf{Key idea:} simulation lets us generate many possible insurance
years and use the simulated aggregate losses to approximate quantities
such as \(E[S]\) and \(\operatorname{Var}(S)\).
\end{quote}
\end{frame}

\section{Properties of the Compound Binomial
Model}\label{properties-of-the-compound-binomial-model}

\begin{frame}{Comparing the Three Frequency Models}
\protect\phantomsection\label{comparing-the-three-frequency-models}
The three models differ mainly in their assumptions about the number of
claims \(N\).

{\def\LTcaptype{none} % do not increment counter
\begin{longtable}[]{@{}
  >{\raggedright\arraybackslash}p{(\linewidth - 4\tabcolsep) * \real{0.3562}}
  >{\raggedright\arraybackslash}p{(\linewidth - 4\tabcolsep) * \real{0.4521}}
  >{\raggedright\arraybackslash}p{(\linewidth - 4\tabcolsep) * \real{0.1918}}@{}}
\toprule\noalign{}
\begin{minipage}[b]{\linewidth}\raggedright
Model
\end{minipage} & \begin{minipage}[b]{\linewidth}\raggedright
Frequency distribution
\end{minipage} & \begin{minipage}[b]{\linewidth}\raggedright
Support of \(N\)
\end{minipage} \\
\midrule\noalign{}
\endhead
Compound Binomial & \(\operatorname{Binomial}(n,p)\) & \(0,\ldots,n\) \\
Compound Poisson & \(\operatorname{Poisson}(\lambda)\) &
\(0,1,2,\ldots\) \\
Compound Negative Binomial & \(\operatorname{NB}(k,p)\) &
\(0,1,2,\ldots\) \\
\bottomrule\noalign{}
\end{longtable}
}

The claim amounts \(X_i\) can have the same severity distribution in all
three models.
\end{frame}

\begin{frame}{Comparing Frequency Variability}
\protect\phantomsection\label{comparing-frequency-variability}
For the binomial model,

\[
E[N]=np,
\qquad
\operatorname{Var}(N)=np(1-p).
\]

For the Poisson model,

\[
E[N]=\lambda,
\qquad
\operatorname{Var}(N)=\lambda.
\]

For the negative binomial model,

\[
E[N]=\frac{kq}{p},
\qquad
\operatorname{Var}(N)=\frac{kq}{p^2}.
\]

Thus, the binomial model has \(\operatorname{Var}(N)<E[N],\)

while the Poisson and negative binomial models have variance equal to or
greater than their mean.
\end{frame}

\begin{frame}{Interpretation of Frequency Variability}
\protect\phantomsection\label{interpretation-of-frequency-variability}
The three models represent different patterns of claim-count
variability.

\textbf{Binomial:} fixed number of potential claim events and limited
frequency variability.

\textbf{Poisson:} no fixed upper limit, with frequency variance equal to
the mean.

\textbf{Negative binomial:} no fixed upper limit and greater frequency
variability than the Poisson model.

The choice of frequency model should reflect the underlying
claim-generation mechanism.
\end{frame}

\begin{frame}{A Fixed Upper Limit on Claims}
\protect\phantomsection\label{a-fixed-upper-limit-on-claims}
A defining feature of the compound binomial model is

\[
N\leq n.
\]

Consequently, the portfolio can contain at most \(n\) claims.

This is natural when there are \(n\) policies and each policy can
produce at most one claim.

In contrast, Poisson and negative binomial models have no finite upper
limit on \(N\).
\end{frame}

\begin{frame}{The Role of the Claim Probability}
\protect\phantomsection\label{the-role-of-the-claim-probability}
For

\[
N\sim\operatorname{Binomial}(n,p),
\]

the expected number of claims is

\[
E[N]=np.
\]

Increasing \(p\) increases the expected frequency.

However,\(\operatorname{Var}(N)=np(1-p),\) so frequency variability is
also affected by \(p\).

The variance is largest when \(p=0.5\).
\end{frame}

\begin{frame}{Aggregate Loss Depends on Frequency and Severity}
\protect\phantomsection\label{aggregate-loss-depends-on-frequency-and-severity}
For the compound binomial model,

\[
E[S]=npE[X].
\]

The expected aggregate loss therefore depends on:

\begin{itemize}
\tightlist
\item
  the number of potential claims \(n\);
\item
  the probability of a claim \(p\);
\item
  the expected claim amount \(E[X]\).
\end{itemize}

Similarly,

\[
\operatorname{Var}(S)
=
E[N]\operatorname{Var}(X)
+
\operatorname{Var}(N)\{E[X]\}^2.
\]

Thus, aggregate risk comes from both \textbf{claim severity} and
\textbf{claim frequency}.
\end{frame}

\begin{frame}{Binomial to Poisson as an Approximation}
\protect\phantomsection\label{binomial-to-poisson-as-an-approximation}
The binomial frequency model approaches a Poisson model when

\[
n\to\infty,
\qquad
p\to0,
\qquad
np\to\lambda.
\]

Then

\[
\operatorname{Binomial}(n,p)
\approx
\operatorname{Poisson}(\lambda).
\]

This provides a useful connection between the compound binomial and
compound Poisson models.

The Poisson model can therefore be viewed as an approximation when there
are many potential claim events with small individual claim
probabilities.
\end{frame}

\begin{frame}{Compound Binomial: Key Results}
\protect\phantomsection\label{compound-binomial-key-results}
For

\[
N\sim\operatorname{Binomial}(n,p),
\]

and

\[
S=\sum_{i=1}^{N}X_i,
\]

with \(N\) independent of the i.i.d. claim amounts \(X_i\),

\[
E[S]=np\,m_1,\qquad \operatorname{Var}(S)=npm_2-np^2m_1^2.
\]

The compound binomial model is particularly appropriate when there is a
\textbf{fixed number of potential claim events}, with each event
producing at most one claim.
\end{frame}

\begin{frame}{Three Compound Models}
\protect\phantomsection\label{three-compound-models}
The main distinction among the three special compound models is their
\textbf{frequency distribution}.

\[
\begin{array}{c|c|c}
\text{Model} & E[N] & \operatorname{Var}(N)\\
\hline
\text{Binomial} & np & np(1-p)\\
\text{Poisson} & \lambda & \lambda\\
\text{Negative Binomial} & \dfrac{kq}{p} & \dfrac{kq}{p^2}
\end{array}
\]

The frequency assumption determines the variability of the aggregate
loss.
\end{frame}

\begin{frame}{Revised request}
\protect\phantomsection\label{revised-request}
\begin{quote}
\textbf{Please add a set of slides on proportional reinsurance using the
example below. Introduce the aggregate claims paid by the insurer and
reinsurer, state the key properties of proportional reinsurance, and
then work through the Pareto--compound Poisson example. Break the
solution into multiple slides and use the Pareto parameterisation
already introduced in the course.}
\end{quote}

I would suggest \textbf{10 slides}: 3 for the concept, 1 for the
question, 5 for the solution, and 1 for the key takeaway. The
calculations are deliberately split so that each slide remains readable.
\end{frame}

\begin{frame}{Proportional Reinsurance}
\protect\phantomsection\label{proportional-reinsurance}
Let \(S\) denote the total aggregate claims from a risk during a given
period.

Let \(S_I\) and \(S_R\) denote the aggregate claims paid by the
\textbf{insurer} and \textbf{reinsurer}, respectively.

Then

\[
\boxed{S=S_I+S_R}.
\]

Under proportional reinsurance, the insurer retains a fixed proportion
\(\alpha\) of every claim.
\end{frame}

\begin{frame}{Proportional Reinsurance}
\protect\phantomsection\label{proportional-reinsurance-1}
For each claim \(X_i\),

\[
Y_i=\alpha X_i,
\qquad
Z_i=(1-\alpha)X_i.
\]

Therefore,

\[
S_I=\sum_{i=1}^N\alpha X_i=\alpha S
\]

and

\[
S_R=\sum_{i=1}^N(1-\alpha)X_i=(1-\alpha)S.
\]
\end{frame}

\begin{frame}{Key Features}
\protect\phantomsection\label{key-features}
\begin{itemize}
\tightlist
\item
  Both the insurer and reinsurer contribute to \textbf{every claim}.
\item
  The insurer retains the same proportion \(\alpha\) of every claim.
\item
  The reinsurer pays the remaining proportion \(1-\alpha\).
\item
  Both parties have \textbf{unlimited liability} unless a separate cap
  is imposed.
\end{itemize}

Thus,

\[
\boxed{S_I=\alpha S,\qquad S_R=(1-\alpha)S}.
\]
\end{frame}

\begin{frame}{Worked Example}
\protect\phantomsection\label{worked-example}
Suppose aggregate claims follow a compound Poisson distribution with

\[
N\sim\operatorname{Poisson}(10)
\]

and individual claim amounts

\[
X\sim Pa(4,1).
\]

The insurer purchases proportional reinsurance with retention
\(\alpha=0.8.\)

Find:

\begin{enumerate}
\tightlist
\item
  the distributions of \(S_I\) and \(S_R\);
\item
  their means and variances;
\item
  compare \(\operatorname{Var}(S_I)+\operatorname{Var}(S_R)\) with
  \(\operatorname{Var}(S)\).
\end{enumerate}
\end{frame}

\begin{frame}{Solution: Set Up the Reinsurance Model}
\protect\phantomsection\label{solution-set-up-the-reinsurance-model}
The original aggregate claim is

\[
S=\sum_{i=1}^{N}X_i.
\]

Under proportional reinsurance with \(\alpha=0.8\),

\[
S_I=0.8S,
\qquad
S_R=0.2S.
\]

Equivalently,

\[
S_I=\sum_{i=1}^{N}0.8X_i,
\qquad
S_R=\sum_{i=1}^{N}0.2X_i.
\]
\end{frame}

\begin{frame}{Solution: Claim Amounts After Reinsurance}
\protect\phantomsection\label{solution-claim-amounts-after-reinsurance}
Define the insurer's payment per claim as

\[
Y_i=0.8X_i.
\]

Define the reinsurer's payment per claim as

\[
Z_i=0.2X_i.
\]

Then

\[
S_I=\sum_{i=1}^{N}Y_i,
\qquad
S_R=\sum_{i=1}^{N}Z_i.
\]

Both are therefore compound Poisson distributions.
\end{frame}

\begin{frame}{Solution: Distribution of the Insurer's Claim}
\protect\phantomsection\label{solution-distribution-of-the-insurers-claim}
Recall the scaling property:

\[
X\sim Pa(\alpha,\lambda)
\quad\Longrightarrow\quad
kX\sim Pa(\alpha,k\lambda).
\]

Since

\[
X_i\sim Pa(4,1)
\]

and

\[
Y_i=0.8X_i,
\]

we obtain

\[
\boxed{Y_i\sim Pa(4,0.8)}.
\]
\end{frame}

\begin{frame}{Solution: Distribution of the Reinsurer's Claim}
\protect\phantomsection\label{solution-distribution-of-the-reinsurers-claim}
Similarly,

\[
Z_i=0.2X_i.
\]

Therefore,

\[
\boxed{Z_i\sim Pa(4,0.2)}.
\]

Hence,

\[
\boxed{
S_I\sim\mathcal{CP}(10,Pa(4,0.8))
}
\]

and

\[
\boxed{
S_R\sim\mathcal{CP}(10,Pa(4,0.2))
}.
\]
\end{frame}

\begin{frame}{Solution: Moments of the Original Severity}
\protect\phantomsection\label{solution-moments-of-the-original-severity}
For

\[
X\sim Pa(4,1),
\]

the mean is

\[
E[X]=\frac{\lambda}{\alpha-1}
=\frac{1}{4-1}
=\frac{1}{3}.
\]

The second moment is

\[
E[X^2]
=
\frac{\Gamma(\alpha-2)\Gamma(3)}
{\Gamma(\alpha)}
\lambda^2.
\]

Thus, \(E[X^2]=\frac{\Gamma(2)\Gamma(3)}{\Gamma(4)}=\frac{1}{3}.\)
\end{frame}

\begin{frame}{Solution: Original Aggregate Moments}
\protect\phantomsection\label{solution-original-aggregate-moments}
For a compound Poisson model,

\[
E[S]=\lambda E[X]
\]

and

\[
\operatorname{Var}(S)=\lambda E[X^2].
\]

Here \(\lambda=10\), so

\[
E[S]
=
10\left(\frac{1}{3}\right)
=
\boxed{\frac{10}{3}}.
\]
\end{frame}

\begin{frame}{Solution: Original Aggregate Variance}
\protect\phantomsection\label{solution-original-aggregate-variance}
Using

\[
\operatorname{Var}(S)=\lambda E[X^2],
\]

we obtain

\[
\operatorname{Var}(S)
=
10\left(\frac{1}{3}\right).
\]

Therefore,

\[
\boxed{\operatorname{Var}(S)=\frac{10}{3}}.
\]
\end{frame}

\begin{frame}{Solution: Insurer's Mean}
\protect\phantomsection\label{solution-insurers-mean}
Since

\[
S_I=0.8S,
\]

we can use the scaling property of expectation:

\[
E[S_I]
=
0.8E[S].
\]

Therefore,

\[
E[S_I]
=
0.8\left(\frac{10}{3}\right)
=
\boxed{\frac{8}{3}}.
\]
\end{frame}

\begin{frame}{Solution: Insurer's Variance}
\protect\phantomsection\label{solution-insurers-variance}
For a constant \(c\),

\[
\operatorname{Var}(cS)=c^2\operatorname{Var}(S).
\]

Therefore,

\[
\operatorname{Var}(S_I)
=
(0.8)^2\operatorname{Var}(S).
\]

Hence,

\[
\operatorname{Var}(S_I)
=
(0.8)^2\left(\frac{10}{3}\right)
=
\boxed{\frac{32}{15}}.
\]
\end{frame}

\begin{frame}{Solution: Reinsurer's Mean}
\protect\phantomsection\label{solution-reinsurers-mean}
Similarly,

\[
S_R=0.2S.
\]

Therefore,

\[
E[S_R]
=
0.2E[S]
=
0.2\left(\frac{10}{3}\right).
\]

Thus,

\[
\boxed{E[S_R]=\frac{2}{3}}.
\]
\end{frame}

\begin{frame}{Solution: Reinsurer's Variance}
\protect\phantomsection\label{solution-reinsurers-variance}
Using the variance scaling property,

\[
\operatorname{Var}(S_R)
=
(0.2)^2\operatorname{Var}(S).
\]

Therefore,

\[
\operatorname{Var}(S_R)
=
(0.2)^2\left(\frac{10}{3}\right)
=
\boxed{\frac{2}{15}}.
\]
\end{frame}

\begin{frame}{Solution: Compare the Variances}
\protect\phantomsection\label{solution-compare-the-variances}
The sum of the two individual variances is

\[
\operatorname{Var}(S_I)+\operatorname{Var}(S_R)
=
\frac{32}{15}+\frac{2}{15}
=
\boxed{\frac{34}{15}}.
\]

The original aggregate variance is

\[
\operatorname{Var}(S)=\frac{10}{3}
=\frac{50}{15}.
\]

Therefore,

\[
\boxed{
\frac{34}{15}<\frac{50}{15}
}.
\]
\end{frame}

\begin{frame}{Solution: Why Are They Different?}
\protect\phantomsection\label{solution-why-are-they-different}
We know that

\[
S=S_I+S_R.
\]

Thus,

\[
\operatorname{Var}(S)
=
\operatorname{Var}(S_I)
+
\operatorname{Var}(S_R)
+
2\operatorname{Cov}(S_I,S_R).
\]

The covariance term is important because

\[
S_I=0.8S,
\qquad
S_R=0.2S.
\]
\end{frame}

\begin{frame}{Solution: The Covariance}
\protect\phantomsection\label{solution-the-covariance}
We have

\[
\operatorname{Cov}(S_I,S_R)
=
\operatorname{Cov}(0.8S,0.2S).
\]

Therefore,

\[
\operatorname{Cov}(S_I,S_R)
=
0.8(0.2)\operatorname{Var}(S).
\]

Hence,

\[
\operatorname{Cov}(S_I,S_R)
=
0.16\left(\frac{10}{3}\right)
=
\boxed{\frac{8}{15}}.
\]

So the variance decomposition becomes
\(\frac{32}{15}+\frac{2}{15}+2\left(\frac{8}{15}\right)=\frac{50}{15}=\frac{10}{3}.\)
\end{frame}

\begin{frame}{Key Takeaways}
\protect\phantomsection\label{key-takeaways}
Under proportional reinsurance,

\[
\boxed{S_I=\alpha S,\qquad S_R=(1-\alpha)S}.
\]

Therefore,

\[
E[S_I]=\alpha E[S],
\qquad
E[S_R]=(1-\alpha)E[S].
\]

and

\[
\operatorname{Var}(S_I)=\alpha^2\operatorname{Var}(S),
\qquad
\operatorname{Var}(S_R)=(1-\alpha)^2\operatorname{Var}(S).
\]

The insurer and reinsurer payments are \textbf{dependent}, so
\(\operatorname{Var}(S_I)+\operatorname{Var}(S_R)\neq\operatorname{Var}(S).\)

This expanded version gives the students more visual space while making
the \textbf{scaling property} and the \textbf{covariance issue} much
easier to follow.
\end{frame}




\end{document}
